\section{Overview of the \gai Benchmark Problem Development Methodology}
\label{section:graph-based-ai-problem-methodology}
% \begin{landscape}
    
    \begin{figure}[h!]
        \begin{center}
            \includesvg[width=\columnwidth]{Figures/SVGs/benchmark_design_process}
            \caption{\gai Benchmark Design Framework}
            \label{figure:framework}
        \end{center}
    \end{figure}

% \end{landscape}

A formal method for designing benchmark problems ensures the transparency, verifiability, and reproducibility of benchmark problems and guides the extension of the benchmark to include emerging \gai problems. 
Figure~\ref{Figure:framework} shows how the process comprises six steps that generate the components of the benchmark specification and the observation apparatus.

The process begins by identifying a real-world use case to ground the problem in real-world datasets and current approaches. 
A more detailed analysis of the problem yields information about the size and shape of the data and formally specifies the everyday tasks associated with that problem in natural language. 
Formalizing the study's phenomena includes identifying what specific interactions of dataset and task are likely to yield interesting results, providing the first significant contribution to the problem's `grey box' approach. 
Considering the phenomena, exploratory analysis of the domain builds the base of knowledge to make initial hypotheses about behavior. 
The hypotheses drive the data collection to study the phenomena of interest, and deviation from those hypotheses indicates interesting behavior worthy of further study. 
The construction of the apparatus involves collecting and generating the datasets that will be used to trigger the conditions the hypotheses predict will stimulate the phenomena of interest. 
It also involves developing any required \ac{API} extensions to facilitate the problem and evaluate the solutions. 
The benchmark user is responsible for all results analysis beyond the return of the efficiency and effectiveness reports by the benchmark. 

The methodology for designing benchmark problems yields a collection of datasets, tasks, and metrics that will combine to provide a `grey box' view of the \ac{SUT} stimulating phenomena of interest according to the defined hypotheses. 
